\documentclass[11pt]{article}
\usepackage{amsmath,amssymb,amsthm,mathtools}\usepackage[margin=1in]{geometry}\usepackage{url}
\newtheorem{theorem}{Theorem}\newtheorem{lemma}[theorem]{Lemma}\newtheorem{proposition}[theorem]{Proposition}
\theoremstyle{remark}\newtheorem{remark}[theorem]{Remark}
\newcommand{\pb}{\overline{p}}
\title{An overpartition congruence modulo $7$:\\ proof of a conjecture of Wei}
\author{Anonymous}\date{12 October 2026}
\begin{document}\maketitle
\begin{abstract}
Let $\pb(n)$ denote the number of overpartitions of $n$. Wei conjectured that $\pb(16n)\equiv0\pmod 7$ whenever $n\equiv3\pmod 8$ and $\left(\frac n7\right)=1$. We prove this. Theta-function $2$-dissections give an exact expression for $\sum\pb(16n)q^n$; reduction modulo $7$ turns it into a holomorphic modular form of weight $18$ on $\Gamma_0(224)$ divided by a power series in $q^{56}$; a twisting argument and Sturm's theorem finish the proof.
\end{abstract}

\section{Introduction}
Overpartitions have generating function
\[\sum_{n\ge0}\pb(n)q^n=\prod_{k\ge1}\frac{1+q^k}{1-q^k}=\frac1{\varphi(-q)},\]
where $\varphi(q)=\sum_{k\in\mathbb Z}q^{k^2}$. Wei~\cite{Wei} proved Ramanujan-type congruences $\pb(11(8n+5))\equiv0\pmod{11}$ and $\pb(13\cdot2^6(8n+7))\equiv0\pmod{13}$ and stated:

\medskip\noindent\textbf{Conjecture} (\cite[Conjecture 4.1]{Wei}). If $n\equiv3\pmod8$ and $\left(\frac n7\right)=1$, then $\pb(2^4n)\equiv0\pmod 7$. Equivalently, $\pb(16(56n+k))\equiv0\pmod7$ for all $n\ge0$ and $k\in\{11,43,51\}$.

\begin{theorem}\label{thm} For every integer $n\ge0$ and $k\in\{11,43,51\}$,
\[\pb(16(56n+k))\equiv0\pmod7.\]
Equivalently, $\pb(16m)\equiv0\pmod7$ for every $m\ge0$ with $m\equiv3\pmod8$ and $\left(\frac m7\right)=1$.\end{theorem}

Ryan et al.~\cite[Theorems 1.3 and 1.5]{RSVZ} already give square-prime, Legendre-restricted families, as well as cube-prime families. Our distinction is the actual target: the arguments here are prime to $7$, whereas their displayed modulo-$7$ families contain that prime as a factor. We claim this specialised congruence, not a new general twisting method. Notation: $q=e^{2\pi iz}$, $\operatorname{Im}z>0$; $E[f](q)=\sum_n f_{2n}q^n$ for $f=\sum f_nq^n$, so $E[f](q^2)=\frac12(f(q)+f(-q))$; and $S_k(q)=\sum_n\pb(kn)q^n$. Negative-index coefficients are zero.

\section{Dissections}
We use three classical identities (Berndt~\cite[Ch.~16, Entry 25]{Berndt}):
\[(\mathrm a)\ \varphi(q)+\varphi(-q)=2\varphi(q^4),\qquad(\mathrm b)\ \varphi(q)\varphi(-q)=\varphi(-q^2)^2,\qquad(\mathrm c)\ \varphi(q)^2+\varphi(-q)^2=2\varphi(q^2)^2.\]
\begin{lemma}\label{lem:dis}
\[S_2=\frac{\varphi(q^2)}{\varphi(-q)^2},\quad S_4=\frac{\varphi(q)^3}{\varphi(-q)^4},\quad S_8=\frac{E[\varphi^7](q)}{\varphi(-q)^8},\quad S_{16}=\frac{E\big[E[\varphi^7]\,\varphi^8\big](q)}{\varphi(-q)^{16}}.\]
\end{lemma}
\begin{proof} By (b), $1/\varphi(-q)=\varphi(q)/\varphi(-q^2)^2$, and $E[\varphi](q)=\varphi(q^2)$ by (a), giving $S_2$. By (b) again, $S_2=\varphi(q^2)\varphi(q)^2/\varphi(-q^2)^4$, and $E[\varphi^2]=\varphi^2$ by (c), giving $S_4$. Writing $S_4=\varphi(q)^7/\varphi(-q^2)^8$ (by (b)) gives $S_8$, and writing $1/\varphi(-q)^8=\varphi(q)^8/\varphi(-q^2)^{16}$ gives $S_{16}$.\end{proof}

\section{Reduction modulo 7}
Reduction modulo $7$ is a ring homomorphism $\mathbb Z[[q]]\to\mathbb F_7[[q]]$ commuting with $E$ and with inversion of series with constant term $1$, and $f(q)^7\equiv f(q^7)$.

\begin{lemma}\label{lem:E8} With $X=\varphi(q)^4$, $Y=\varphi(-q)^4$: $E[\varphi^8]=8X^2-8XY+Y^2$.\end{lemma}
\begin{proof} Put $a=\varphi(t)^2$, $b=\varphi(-t)^2$ with $t^2=q$. By (c) and (b), $a+b=2\varphi(q)^2$ and $ab=\varphi(-q)^4$, so $a^2+b^2=4X-2Y$ and $a^4+b^4=(4X-2Y)^2-2Y^2$. Then $E[\varphi^8](q)=\frac12(a^4+b^4)$.\end{proof}

Since $\varphi^7\equiv\varphi(q^7)$, $E[\varphi^7]\equiv\varphi(q^{14})$ and 
\[E[E[\varphi^7]\varphi^8]\equiv\varphi(q^7)E[\varphi^8]\equiv\varphi(q^7)(X^2-XY+Y^2)\qquad(\text{as }8\equiv1).\] Also $1/\varphi(-q)^{16}\equiv\varphi(-q)^5/\varphi(-q^7)^3$, and iterating (b),
\[\frac1{\varphi(-q^7)^3}=\frac{\varphi(q^7)^3\varphi(q^{14})^6\varphi(q^{28})^{12}}{\varphi(-q^{56})^{24}}.\] Multiplying numerator and denominator by $\varphi(q^{56})$:
\begin{equation}\label{eq:red} S_{16}(q)\equiv\frac{\widetilde H(q)}{G(q^{56})}\pmod 7,\end{equation}
where
\[\widetilde H=(X^2-XY+Y^2)\,\varphi(-q)^5\varphi(q^7)^4\varphi(q^{14})^6\varphi(q^{28})^{12}\varphi(q^{56}),\qquad G=\varphi(-q)^{24}\varphi(q).\]
The extra factor $\varphi(q^{56})$ makes the numerator's weight integral while preserving the denominator's support on multiples of $56$ and its constant term $1$.
Write $\widetilde H=\sum K(m)q^m$. Since $1/G(q^{56})$ is a series in $q^{56}$, $\pb(16n)\equiv\sum_{J=0}^{\lfloor n/56\rfloor}e_JK(n-56J)$ with $n-56J\equiv n\pmod{56}$. Hence Theorem~\ref{thm} follows from:

\begin{proposition}\label{prop:K} $K(m)\equiv0\pmod 7$ for all $m\equiv3\pmod8$ with $\left(\frac m7\right)=1$.\end{proposition}

\section{Modularity and the Sturm bound}
With $\varphi(q^d)=\eta(2d\tau)^5/(\eta(d\tau)^2\eta(4d\tau)^2)$ and $\varphi(-q^d)=\eta(d\tau)^2/\eta(2d\tau)$, each of the three monomials $X^2Z$, $XYZ$, $Y^2Z$ of $\widetilde H$ (with $Z$ the product of the remaining theta factors) is an eta-quotient of weight $18$ and level $224$. By the Gordon--Hughes--Newman criterion and Ligozat's cusp-order formula~\cite[Thms.~1.64, 1.65]{Ono}, each lies in $M_{18}(\Gamma_0(224),\chi)$ with $\chi=\left(\frac{14}{\cdot}\right)$ (we checked $\sum\delta r_\delta\equiv\sum(224/\delta)r_\delta\equiv0\pmod{24}$, the character, and nonnegative order at every cusp in exact arithmetic). Hence $\widetilde H\in M_{18}(\Gamma_0(224),\chi)\cap\mathbb Z[[q]]$.

\begin{lemma}[Twisting at a sufficient level]\label{lem:tw} Let $f=\sum a(m)q^m\in M_k(\Gamma_0(N),\chi)$, $k\in\mathbb Z$, let $C$ be the conductor of $\chi$, and let $\psi$ be any Dirichlet character modulo $M$. Then, for $L=\operatorname{lcm}(N,M^2,MC)$,
\[f\otimes\psi=\sum\psi(m)a(m)q^m\in M_k(\Gamma_0(L),\chi\psi^2).\]
No minimality of $L$ or novelty of this twisting principle is asserted.\end{lemma}
\begin{proof} With $T_j=\begin{psmallmatrix}1&j/M\\0&1\end{psmallmatrix}$ and $c_j=M^{-1}\sum_{r\bmod M}\psi(r)\zeta_M^{-rj}$ we have $f\otimes\psi=\sum_{j\bmod M}c_j f|T_j$. For $\gamma=\begin{psmallmatrix}a&b\\c&d\end{psmallmatrix}\in\Gamma_0(L)$, multiplication gives
\[T_j\gamma T_{-d^2j}=\begin{pmatrix}a+cj/M&b-bcdj/M-cd^2j^2/M^2\\c&d-cd^2j/M\end{pmatrix}.\]
Its entries are integral since $M^2\mid c$, its lower-left entry is divisible by $N$, and its lower-right entry is $d$ modulo $C$ since $MC\mid c$. Thus $(f|T_j)|\gamma=\chi(d)f|T_{d^2j}$. Reindexing gives $c_{d^{-2}j}=\psi(d)^2c_j$, proving the law, even for imprimitive $\psi$. For any rational matrix $\alpha$ of positive determinant, choose $\sigma\in\mathrm{SL}_2(\mathbb Z)$ with $\sigma\infty=\alpha\infty$. Then $\sigma^{-1}\alpha$ is upper triangular with positive dilation ratio; it preserves nonnegative exponents in a holomorphic cusp expansion. Apply this to each translated cusp matrix $\alpha=T_j\gamma$.\end{proof}

The indicator of $\{m\equiv3\ (8),\ (\frac m7)=1\}$ is a linear combination of real characters modulo $56$:
\[\frac14\sum_{\psi\bmod8}\psi(3)\psi(m)\cdot\frac12\big(\chi_0^{(7)}(m)+(\tfrac m7)\big).\]
Here $C=56$ and $L=\operatorname{lcm}(224,56^2,56^2)=3136$. All twisting characters square to the principal character on units, so
\[g:=\sum_{m\equiv3(8),\,(m/7)=1}K(m)q^m\in M_{18}(\Gamma_0(3136),\chi)\cap\mathbb Z[[q]].\]
By Sturm's theorem with character~\cite{Sturm,Ono}, the sufficient inclusive bound is
\[B=\frac{18}{12}\cdot3136\cdot\frac32\cdot\frac87=8064.\]

\begin{proposition}[computation]\label{prop:comp} For all $m\le8064$ with $m\equiv3\pmod8$ and $\left(\frac m7\right)=1$ (there are $432$ such $m$), $K(m)\equiv0\pmod7$. There are $370$ nonzero complementary controls and $K(0)=1$.\end{proposition}
This proves Proposition~\ref{prop:K} and hence Theorem~\ref{thm}.

\section{Computations}
All proof checks are exact. From the collection root run \texttt{python verify.py mod7}. The identity checker compares the dissections with a separate overpartition recurrence; the modularity checker verifies the eta conditions; and both FLINT and a standard-library implementation compute the inclusive bound $8064$, obtaining the $432$ zero coefficients and $370$ nonzero complementary controls of Proposition~\ref{prop:comp}. The original larger-bound calculation is historical corroboration, not required by this proof. These are reproducible computational certificates, not formal proof-assistant verification.

For a compact modularity audit, the exponent vectors of $X^2Z,XYZ,Y^2Z$, in the order $\delta=(1,2,4,7,14,28,56,112,224)$, are
\[(-6,35,-16,-8,8,-2,46,-19,-2),\quad
(10,11,-8,-8,8,-2,46,-19,-2),\]
\[(26,-13,0,-8,8,-2,46,-19,-2).\]
Their orders at cusp denominators $(1,2,4,7,8,14,16,28,32,56,112,224)$ are respectively
\[(70,120,6,10,12,72,2,42,0,84,14,0),\]
\[(126,64,6,18,12,64,2,42,0,84,14,0),\]
\[(182,8,6,26,12,56,2,42,0,84,14,0).\]
These are width-normalised cusp orders, so their nonnegativity proves holomorphy at every cusp.

\section*{Collection and verification}
This paper is part of the revised collection, version 1.1.0-candidate, DOI \url{https://doi.org/10.5281/zenodo.23319157}. Companion papers and exact verification code are archived together. The statements are unrefereed computer-assisted proof candidates; no formal verification or unaffiliated reproduction is asserted.

\begin{thebibliography}{6}
\bibitem{Berndt} B.~C.~Berndt, \emph{Ramanujan's Notebooks, Part III}, Springer, New York, 1991.
\bibitem{Ono} K.~Ono, \emph{The Web of Modularity: Arithmetic of the Coefficients of Modular Forms and $q$-series}, CBMS Regional Conference Series in Mathematics 102, AMS, 2004.
\bibitem{RSVZ} N.~C.~Ryan, N.~Sirolli, J.~C.~Villegas-Morales and Q.-Y.~Zheng, Explicit families of congruences for the overpartition function, \emph{Ramanujan J.} (2024); arXiv:2309.01792.
\bibitem{Sturm} J.~Sturm, On the congruence of modular forms, in \emph{Number Theory (New York, 1984--1985)}, Lecture Notes in Math. 1240, Springer, 1987, 275--280.
\bibitem{Wei} X.~Wei, New Ramanujan-type congruences for overpartitions modulo 11 and 13, arXiv:2603.08510 (2026).
\end{thebibliography}
\end{document}
